\scanpage{51}
Conditions pour une partition admissible

\textcircled{a} feuilles connexes

\textcircled{b} feuilles non-singulières (contextes d'espaces \inserted{\unsure{top. ou}} L.P.M. ou S.A.P.M.\ldots)

\textcircled{c} équisingularité

\uline{NB} a) n'est pas stable en s'induisant sur un ouvert de $X$ (tandis que la condition d'admissibilité l'est, ainsi que b) et c)). Une bonne condition stable est

\inserted{Les $F_i$ sont \inserted{connexes et} loc. connexes, et}

(C) $\forall$ ouvert $U$ de $X$, la famille de \deleted{toutes} les composantes connexes de tous les $U\cap F_i$ est une partition admissible de $U$,
\inserted{a) $\forall i$, la famille des comp. conn. de $U\cap F_i$ est loc. finie sur $U$ et}
\deleted{\unsure{alors}} \inserted{b)} si $F_{i\alpha}$ $(F_{j\beta})$ une composante connexe de $U\cap F_i$ \inserted{$(U\cap F_j)$} \deleted{$\unsure{U\cap F_i}$}
\[
F_{j\beta}\cap\overline F_{i\alpha}\neq\varnothing
\Longrightarrow F_{j\beta}\subset\overline F_{i\alpha}
\]
(ou encore, $\overline F_{i\alpha}\cap F_j\cap U$ est une partie ouverte de $F_j\cap U$ \deleted{$\unsure{F_{j\beta}}$}, ou encore de $F_j$)

Soit $(F_i)$ partition admissible.
\[
F_i\text{ ouvert dans }X\Longleftrightarrow\overline F_i\text{ maximal.}
\]
$F_i$ ouvert \inserted{de $X$, et} $\overline F_i\subsetneq\overline F_j$
$\Rightarrow F_i$ ouvert \inserted{$\neq\varnothing$} de $\overline F_j$ et $\overline F_i$ rare dans $\overline F_j$, absurde

donc $F_i$ ouvert dans $X\longrightarrow F_i$ maximal

Soit $F_i$ maximal, alors
\[
X-F_i=\bigcup_{j\neq i}\overline F_j\quad(\text{fermé})
\]
car \unsure{$F_j$} si $j\neq i$, $\overline F_j\cap F_i=\varnothing$

$F_i$ ouvert $\Longleftrightarrow X-F_i\;(=\bigcup_{j\neq i}F_j)$ fermé
$\Longleftrightarrow\forall j\neq i,\;\overline F_j\subset X-F_i$

i.e. $\forall j\neq i$, \deleted{$\unsure{F_j=\overline F_j}$} $k$ tel que $\overline F_k\subsetneq\overline F_j$, on a $k\neq i$

$\Longleftrightarrow i$ maximal.

\margin{52}
